\begin{Note}[{\it Alternative proof of Lemma} \ref{lem2}] \label{Quest4.3}
In the right-hand side of Lemma \ref{lem2}, the quantity
$\binom{k_1+s_1}{k_1}\cdots \binom{k_d+s_d}{k_d}$ appearing there can also be interpreted as
the pairing
$\lp w_{(k_1,\dots,k_d;0)}\sha X^s, w_{(k_1+s_1,\dots,k_d+s_d;0)} \rp$.
Therefore, the assertion of this lemma is equivalent to the identity
\begin{equation}
\sum_{\btau\in\N_0^d} 
\lp\SS^{(\btau;0)},w_{(\bkappa+\bs;0)} \rp \cdot \lp\MM^{(\btau;0)}, w_{(\bkappa;0)}\!\cdot\!\! X^s\rp
=(-1)^s 
\lp w_{(\bkappa;0)}\sha X^s, w_{(\bkappa+\bs;0)} \rp
\label{alterLemma}
\end{equation}
for $\bkappa=(k_1,\dots,k_d)\in\N_0^d$, $\bs=(s_1,\dots,s_d)\in\Z^d$ 
satisfying $s=s_1+\cdots+s_d\ge 0$ and $\bkappa+\bs\in\N_0^d$.
We now give an alternative proof for it using the Magnus/demi-shuffle duality.
First, by Corollary \ref{coro}, we have 
$w_{(\bkappa;0)}=\sum_\bbr\la \MM^{(\bbr)},w_{(\bkappa;0)}\ra \SS^{(\bbr)}$ 
and $w_{(\bkappa+\bs;0)}=\sum_\bt\la \SS^{(\bt)},w_{(\bkappa+\bs;0)}\ra \MM^{(\bt)}$
so that the right-hand side of \eqref{alterLemma} can be written as
\begin{align}
\label{RHS4.2}
&(-1)^s 
\lp w_{(\bkappa;0)}\sha X^s, w_{(\bkappa+\bs;0)} \rp \\
&=(-1)^s \sum_{\bbr,\bt\in\N_0^{(\infty)}}
\lp \SS^{(\bbr)}\sha X^s, \MM^{(\bt)} \rp
\lp \MM^{(\bbr)} ,w_{(\bkappa;0)} \rp
\lp \SS^{(\bt)} ,w_{(\bkappa+\bs;0)}\rp
\notag \\
&=(-1)^s \sum_{\brho\in\N_0^d} 
\lp \MM^{(\brho;0)} ,w_{(\bkappa;0)}\rp
\lp  \SS^{(\brho;s)}, w_{(\bkappa+\bs;0)} \rp.
\notag
\end{align}
Here in the second equality, we use the fact that
$\lp  \MM^{(\bbr)}, w_{(\bkappa;0)} \rp$ 
survives only if $\bbr=(\brho;0)\in\N_0^{(\infty)}$ for some $\brho\in\N_0^d$
and then apply the duality (Theorem \ref{orthogonality})
to $\lp \SS^{(\bbr)}\sha X^s, \MM^{(\bt)} \rp$
with $\SS^{(\brho;0)}\sha X^s=\SS^{(\brho;s)}$ 
(cf. Definitions \ref{defMagnus} and \ref{demi-shuffle}).

On the other hand, in the left-hand side of  \eqref{alterLemma},
one observes that nontrivial terms of the summation arise
only from those $\btau=(\tau_1,\dots,\tau_d)\in\N_0^d$ with
$\sum_{i=1}^d \tau_i=s+\sum_{i=1}^d \kappa_i$ (a constant). 
But the last binomial factor in \eqref{MagnusBino} for
$ \lp\MM^{(\btau;0)}, w_{(\bkappa;0)}\!\cdot\!\! X^s\rp
=\magnus{(\tau_1,\dots,\tau_d;0)}{(\kappa_1,\dots,\kappa_d;s)}
$
equals $\binom{\tau_d}{s}$, which is non-zero only if
$\tau_d\ge s$.
Therefore, 
the summation~$\sum_\btau$ may be replaced by 
$\sum_\brho$ with $\brho=\btau-(\mathbf{0},s)$
in $\N_0^d$ (where $\mathbf{0}\in\N_0^{d-1}$ is the zero vector).
Thus,
the left-hand side of \eqref{alterLemma} can be written as
\begin{align}
\label{LHS4.2}
&\sum_{\btau\in\N_0^d} 
\lp\SS^{(\btau;0)},w_{(\bkappa+\bs;0)} \rp \cdot 
\lp\MM^{(\btau;0)}, w_{(\bkappa;0)}\!\cdot\!\! X^s\rp \\
&=
\sum_{\brho\in\N_0^d} 
\lp\SS^{(\brho+(\mathbf{0},s);0)},w_{(\bkappa+\bs;0)} \rp \cdot 
\lp\MM^{(\brho+(\mathbf{0},s);0)}, w_{(\bkappa;s)} \rp.
\notag
\end{align}
Comparing summands of the above \eqref{RHS4.2} and \eqref{LHS4.2} for 
individual $\brho\in\N_0^d$ in view of coefficients of 
monomial expansions of demi-shuffle/Magnus polynomials 
(Lemma~\ref{dS-mono} and~\eqref{Mag-Mono}), we reduce 
the formula \eqref{alterLemma} to the following elementary identity
for $\bkappa=(k_i)$, $\brho=(r_i)\in\N_0^d$ and $\bs=(s_i)\in\Z^d$
satisfying
$\sum_{i=1}^d k_i=\sum_{i=1}^d r_i$, $\bs+\bkappa\in \N_0^d$
and $s:=\sum_{i=1}^d s_i\ge 0$:
\begin{equation}
\binom{(\bkappa+\bs;0)}{(\brho;s)}
\magnus{(\brho;0)}{(\bkappa;0)}
=
(-1)^s
\binom{(\bkappa+\bs;0)}{(\brho+(\mathbf{0},s);0)}
\magnus{(\brho+(\mathbf{0},s);0)}{(\bkappa,s)}.
\end{equation} 
This is an immediate consequence of definitions
of these symbols $\{^*_* \}$, $(^*_*)$.
(Observe that only difference between the corresponding symbols
occurs from the last binomial coefficient in  \eqref{ArrayBino}
and \eqref{MagnusBino}.)
\qed
\end{Note}
